% Part:sets-functions-relations
% Chapter: size-of-sets
% Section: enumerations

\documentclass[../../../include/open-logic-section]{subfiles}

\begin{document}

\olfileid{sfr}{siz}{enm}

\olsection{પરિગણનાઓ અને \usetoken{S}{enumerable} ગણો}

\begin{editorial}
આ વિભાગ ગણોની પરિગણનાઓને $\PosInt$થી આવતા વ્યાપ્ત
વિધેયો તરીકે વ્યાખ્યાયિત કરીને તેમની ચર્ચા કરે છે.
ઓછી ગાણિતિક પૃષ્ઠભૂમિ ધરાવતા વાચકો માટે વાત ધીમે ધીમે
આગળ વધે છે. \olref[enm-alt]{sec}માં વધુ સંક્ષિપ્ત
વૈકલ્પિક સ્વરૂપ આપેલું છે. તેમાં પરિગણનાઓની જુદી વ્યાખ્યા
છે: $\Nat$ (અથવા તેના આરંભખંડ) સાથેનાં એક-એક
વ્યાપ્ત વિધેયો.
\end{editorial}

\begin{explain}
આપણે ગણોના !!{element}sની યાદી આપીને તેમના ઉદાહરણો
જોઈ ચૂક્યા છીએ. હવે વધુ સામાન્ય રીતે ચર્ચા કરીએ કે
કોઈ ગણ અનંત હોય તો પણ તેના !!{element}sની યાદી
કેવી રીતે અને ક્યારે બનાવી શકાય.
\end{explain}

\begin{defn}[પરિગણના, અનૌપચારિક રીતે]
અનૌપચારિક રીતે, ગણ~$A$ની \emph{પરિગણના} એ
$A$ના !!{element}sની એવી યાદી (કદાચ અનંત) છે
કે $A$નો દરેક !!{element} યાદીમાં કોઈ સાન્ત સ્થાને
આવે. જો $A$ને પરિગણના હોય, તો $A$ને
\emph{!!{enumerable}} કહે છે.
\end{defn}

\begin{explain}
પરિગણનાઓ વિશે કેટલાક મુદ્દા:
\begin{enumerate}
\item આપણે એવી જ યાદીઓને પરિગણના ગણીએ છીએ જેમની
  શરૂઆત હોય અને જેમાં પહેલા સિવાયના દરેક !!{element}ની
  તરત પહેલાં એક જ !!{element} હોય. બીજા શબ્દોમાં,
  યાદીના પ્રથમ !!{element} અને બીજા કોઈ પણ !!{element}
  વચ્ચે માત્ર સાન્ત સંખ્યામાં ઘટકો હોય. ખાસ કરીને,
  આનો અર્થ એ છે કે પરિગણનાના દરેક !!{element}નું
  સ્થાન સાન્ત છે: પ્રથમ !!{element}નું સ્થાન~$1$,
  બીજાનું સ્થાન~$2$, વગેરે.
\item એક જ ગણ~$A$ની જુદી પરિગણનાઓ હોઈ શકે,
  જેમાં !!{element}s આવવાનો ક્રમ જુદો હોય:
  $4$, $1$, $25$, $16$,~$9$ એ પ્રથમ પાંચ
  વર્ગસંખ્યાઓના (ગણની) એટલી જ સારી પરિગણના છે
  જેટલી $1$, $4$, $9$, $16$,~$25$ છે.
\item પુનરાવર્તનવાળી પરિગણનાઓ પણ પરિગણનાઓ છે:
  $1$, $1$, $2$, $2$, $3$, $3$,~\dots{}
  એ જ ગણની પરિગણના કરે છે જેની $1$, $2$,
  $3$,~\dots{} પરિગણના કરે છે.
\item પરિગણના નિર્દિષ્ટ કરતી વખતે ક્રમ અને પુનરાવર્તનનું
  મહત્ત્વ \emph{છે}: ધન પૂર્ણાંકોની પરિગણના
  $1$, $2$, $3$, $1$, \dots{}થી શરૂ કરી શકીએ,
  પરંતુ પ્રચલિત રીતે $1$, $2$, $3$, $4$,~\dots
  મુજબ પરિગણના કરીએ ત્યારે રીત વધુ સરળતાથી દેખાય છે.
\item પરિગણનાની શરૂઆત હોવી જોઈએ: \dots, $3$, $2$, $1$
  એ ધન પૂર્ણાંકોની પરિગણના નથી, કારણ કે તેને પ્રથમ
  !!{element} નથી. અનૌપચારિક વ્યાખ્યામાંથી આ વાત
  કેવી રીતે મળે છે તે જોવા પોતાને પૂછો: “યાદીમાં
  સંખ્યા ૭૬ કયા સ્થાને આવે છે?”
\item નીચેની યાદી ધન પૂર્ણાંકોની પરિગણના નથી:
  $1$, $3$, $5$, \dots, $2$, $4$, $6$, \dots\@
  મુશ્કેલી એ છે કે યુગ્મ સંખ્યાઓ સાન્ત સ્થાનોને બદલે
  $\infty + 1$, $\infty + 2$, $\infty +
  3$ સ્થાનો પર આવે છે.
\item ખાલી ગણ ગણનીય છે: ખાલી યાદી તેની પરિગણના કરે છે!{}
\end{enumerate}
\end{explain}

\begin{prop}
જો $A$ને પરિગણના હોય, તો તેને પુનરાવર્તન વિનાની
પરિગણના પણ હોય છે.
\end{prop}

\begin{proof}
ધારો કે $A$ને $x_1$, $x_2$, \dots{} એવી પરિગણના
છે, જેમાં દરેક $x_i$ એ $A$નો !!{element} છે.
વારંવાર આવતા !!{element}s દૂર કરીને પરિગણનામાંથી
પુનરાવર્તન કાઢી શકીએ. દાખલા તરીકે, જો $x_i$ એ
$A$નો એવો !!a{element} હોય જે $x_1$, \dots,
$x_{i-1}$માં ન હોય, તો તેને યાદીમાં રાખીએ;
અને જો $x_i$ પહેલેથી જ $x_1$, \dots,~$x_{i-1}$માં
આવતો હોય, તો તેને યાદીમાંથી દૂર કરીએ. આમ પરિગણનાને
નવી પરિગણનામાં ફેરવી શકીએ.
\end{proof}

છેલ્લી દલીલ બતાવે છે કે પરિગણનાઓ અને !!{enumerable}
ગણોને બરાબર સમજવા તથા તેમના વિશે પરિણામો સાબિત કરવા
વધુ ચોક્કસ વ્યાખ્યા જરૂરી છે. નીચેની વ્યાખ્યા તે પૂરી પાડે છે.

\begin{defn}[પરિગણના, ઔપચારિક રીતે]
ગણ $A \neq \emptyset$ની \emph{પરિગણના} એ કોઈ પણ
!!{surjective} વિધેય $f \colon \PosInt \to A$ છે.
\end{defn}

\begin{explain}
ઔપચારિક વ્યાખ્યા અને કદાચ અનંત યાદી વાપરતી અનૌપચારિક
વ્યાખ્યા સમકક્ષ છે તેની ખાતરી કરીએ. પ્રથમ,
$\PosInt$થી ગણ~$A$માંનું કોઈ પણ !!{surjective}
વિધેય $A$ની પરિગણના કરે છે. આવું વિધેય ઉપરની
અનૌપચારિક વ્યાખ્યા પ્રમાણેની પરિગણના નક્કી કરે છે:
યાદી $f(1)$, $f(2)$, $f(3)$, \dots.
$f$ એ !!{surjective} હોવાથી $A$નો દરેક !!{element}
કોઈક~$n \in \PosInt$ માટે $f(n)$ની કિંમત
ચોક્કસ હોય છે. તેથી $A$નો દરેક !!{element} યાદીમાં
કોઈ સાન્ત સ્થાને આવે છે. વિધેય !!{injective}
ન પણ હોય, તેથી યાદીમાં પુનરાવર્તન હોઈ શકે;
પણ ઉપર નોંધ્યું તેમ તે સ્વીકાર્ય છે.

બીજી તરફ, $A$ના બધા !!{element}sની પરિગણના
કરતી યાદી આપેલી હોય, તો !!a{surjective} વિધેય
$f\colon \PosInt \to A$ આ રીતે વ્યાખ્યાયિત કરી શકીએ:
$f(n)$ એ યાદીનો $n$મો !!{element} હોય; અથવા
$n$મો !!{element} ન હોય તો યાદીનો છેલ્લો
!!{element} હોય. માત્ર $A$ ખાલી હોય અને તેથી
યાદી ખાલી હોય તેવા કિસ્સામાં આ રીત !!a{surjective}
વિધેય આપતી નથી. એટલે દરેક અખાલી યાદી
!!a{surjective} વિધેય $f\colon \PosInt \to A$
નક્કી કરે છે.
\end{explain}

\begin{defn}
  \ollabel{defn:enumerable}
ગણ~$A$ એ !!{enumerable} હોય જો અને માત્ર જો
તે ખાલી હોય અથવા તેને પરિગણના હોય.
\end{defn}

\begin{ex}
ધન પૂર્ણાંકો ($\PosInt$)ની પરિગણના કરતું વિધેય
એ $f(n) = n$ દ્વારા આપેલું તદેવ વિધેય જ છે.
પ્રાકૃતિક સંખ્યાઓ $\Nat$ની પરિગણના કરતું વિધેય
$g(n) = n - 1$ છે.
\end{ex}

\begin{ex}
વિધેયો $f\colon \PosInt \to \PosInt$ અને
$g \colon \PosInt \to
\PosInt$ને
\begin{align*}
f(n) & = 2n \text{ અને}\\
g(n) & = 2n - 1
\end{align*}
દ્વારા આપીએ, તો તેઓ અનુક્રમે યુગ્મ ધન પૂર્ણાંકો અને
અયુગ્મ ધન પૂર્ણાંકોની પરિગણના કરે છે. જોકે, કોઈ પણ
વિધેય $\PosInt$ની પરિગણના નથી, કારણ કે બંનેમાંથી
એકેય !!{surjective} નથી.
\end{ex}

\begin{prob}
ધન વર્ગસંખ્યાઓ $1$, $4$, $9$, $16$, \dots{}ની
પરિગણના વ્યાખ્યાયિત કરો.
\end{prob}

\begin{ex}
વિધેય $f(n) = (-1)^{n} \lceil \frac{(n-1)}{2}\rceil$
પૂર્ણાંકોના ગણ~$\Int$ની પરિગણના કરે છે. અહીં
$\lceil x \rceil$ એ \emph{ઊર્ધ્વ પૂર્ણાંક} વિધેય
દર્શાવે છે, જે $x$ને ઉપરની નજીકની પૂર્ણાંક કિંમત
સુધી લઈ જાય છે. જુઓ કે $f$ કેવી રીતે ધન અને ઋણ
પૂર્ણાંકો વચ્ચે આગળપાછળ “કૂદીને” $\Int$ની કિંમતો
ઉત્પન્ન કરે છે:
\[
\begin{array}{c c c c c c c c}
f(1) & f(2) & f(3) & f(4) & f(5) & f(6) & f(7) & \dots \\ \\
- \lceil \tfrac{0}{2} \rceil & \lceil \tfrac{1}{2}\rceil & - \lceil \tfrac{2}{2} \rceil & \lceil \tfrac{3}{2} \rceil & - \lceil \tfrac{4}{2} \rceil  & \lceil \tfrac{5}{2}
\rceil & - \lceil \tfrac{6}{2} \rceil & \dots \\ \\
0 & 1 & -1 & 2 & -2 & 3 & -3 & \dots
\end{array}
\]
\sourcecorrection{OLSIZ-001}{મૂળની \texttt{enumerability.tex}, પંક્તિઓ
140--153માં કોષ્ટકનું મથાળું અને સૂત્રવાળી હાર \(f(7)=-3\) દર્શાવે છે,
પરંતુ છેલ્લી મૂલ્ય-હારમાં \(-3\) છૂટી ગયું છે. અહીં તે કોષ પૂરો કર્યો છે.}
નીચે મુજબ કિસ્સા પાડીને $f$ વ્યાખ્યાયિત થયું છે
એમ પણ વિચારી શકો:
\[
f(n) = \begin{cases}
  0 & \text{જો $n = 1$ હોય}\\
  n/2 & \text{જો $n$ યુગ્મ હોય}\\
  -(n-1)/2 & \text{જો $n$ અયુગ્મ હોય અને $>1$ હોય}
  \end{cases}
\]
\end{ex}

\begin{prob}
બતાવો કે જો $A$ અને $B$ એ !!{enumerable} હોય,
તો $A \cup B$ પણ છે. તે માટે ધારો કે
!!{surjective} વિધેયો $f\colon \PosInt \to
A$ અને $g\colon \PosInt \to B$ છે; પછી
!!a{surjective} વિધેય~$h\colon \PosInt \to A \cup B$
વ્યાખ્યાયિત કરો અને તે !!{surjective} છે તે સાબિત કરો.
$A$ અથવા~$B = \emptyset$ હોય તેવા કિસ્સાઓ પણ વિચારો.
\end{prob}

\begin{prob}
બતાવો કે જો $B \subseteq A$ હોય અને $A$ એ
!!{enumerable} હોય, તો $B$ પણ છે. તે માટે ધારો કે
!!a{surjective} વિધેય $f\colon \PosInt \to
A$ છે. !!a{surjective} વિધેય~$g\colon \PosInt \to B$
વ્યાખ્યાયિત કરો અને તે !!{surjective} છે તે સાબિત કરો.
જો $B = \emptyset$ હોય તો શું થાય?
\end{prob}

\begin{prob}
$n$ પર ગાણિતિક અનુમાનથી બતાવો કે જો $A_1$, $A_2$,
\dots, $A_n$ બધા !!{enumerable} હોય, તો
$A_1 \cup \dots \cup A_n$ પણ છે.
જો બે ગણો $A$ અને~$B$ એ !!{enumerable} હોય,
તો $A \cup
B$ પણ છે તે હકીકત તમે ધારી શકો છો.
\end{prob}

ગણના !!{element}sની યાદી બનાવતાં $1$લા !!{element}થી
ગણતરી શરૂ કરવી કદાચ વધુ સ્વાભાવિક લાગે, છતાં ગણિતજ્ઞો
વસ્તુઓ ગણવા માટે પ્રાકૃતિક સંખ્યાઓ~$\Nat$ વાપરવાનું
પસંદ કરે છે. તેઓ યાદીના $0$મા, $1$લા, $2$જા વગેરે
!!{element}sની વાત કરે છે. તેને અનુરૂપ,
$\Nat$થી $A$માંનું !!a{surjective} વિધેય લઈને
પરિગણના વ્યાખ્યાયિત કરી શકીએ. અલબત્ત, બંને વ્યાખ્યાઓ
સમકક્ષ છે.

\begin{prop}\ollabel{prop:enum-shift}
!!a{surjection} $f\colon \PosInt \to A$ હોય
જો અને માત્ર જો !!a{surjection} $g\colon \Nat \to A$ હોય.
\end{prop}

\begin{proof}
!!a{surjection} $f\colon \PosInt \to A$ આપેલું હોય,
તો બધા $n \in \Nat$ માટે $g(n) =
f(n+1)$ વ્યાખ્યાયિત કરી શકીએ.
$g\colon \Nat
\to A$ એ !!{surjective} છે તે સરળતાથી જોઈ શકાય.
ઊલટું, !!a{surjection} $g\colon
\Nat \to A$ આપેલું હોય, તો $f(n) = g(n-1)$ વ્યાખ્યાયિત કરો.
\end{proof}

આ પરથી નીચેનું પરિણામ મળે છે:

\begin{cor}\ollabel{cor:enum-nat}
ગણ $A$ એ !!{enumerable} હોય જો અને માત્ર જો
તે ખાલી હોય અથવા !!a{surjective} વિધેય
$f\colon \Nat \to A$ હોય.
\end{cor}

ઉપર ચર્ચા કરી કે ગણ~$A$ના !!{element}sની યાદીને
પુનરાવર્તન વિનાની યાદીમાં ફેરવી શકાય. પરિગણનાઓ માટે
પણ આ સાચું છે, પરંતુ તેને ચોક્કસ રીતે વ્યક્ત કરવું
અને સાબિત કરવું થોડું વધુ અઘરું છે. કોઈ પણ વિધેય
$f\colon \PosInt \to A$ બધા $n \in
\PosInt$ માટે વ્યાખ્યાયિત હોવું જોઈએ. જો $A$માં
માત્ર સાન્ત સંખ્યામાં !!{element}s હોય, તો
$\PosInt$ના અનંત સંખ્યામાં !!{element}s પર
વ્યાખ્યાયિત એવું વિધેય હોઈ ન શકે જે $A$ના બધા
!!{element}sને કિંમતો તરીકે લે, છતાં કોઈ કિંમત
બે વાર ન લે. આ કિસ્સામાં, એટલે કે પુનરાવર્તન
વિનાની યાદી સાન્ત હોય ત્યારે, $f$ માટે જુદો પ્રદેશ
પસંદ કરવો પડે: માત્ર સાન્ત સંખ્યામાં !!{element}s
ધરાવતો પ્રદેશ. પુનરાવર્તન ન હોવાનો અર્થ એ છે
કે $f$ એ !!{injective} હોવું જોઈએ. તે
!!{surjective} પણ હોવાથી, આપણે કોઈ સાન્ત ગણ
$\{1, \dots, n\}$ અથવા $\PosInt$ અને~$A$
વચ્ચે !!a{bijection} શોધીએ છીએ.

\begin{prop}\ollabel{prop:enum-bij}
જો $f\colon \PosInt \to A$ એ !!{surjective}
હોય (એટલે કે $A$ની પરિગણના હોય), તો
!!a{bijection} $g\colon Z \to A$ હોય છે, જ્યાં
$Z$ એ $\PosInt$ અથવા કોઈક~$n \in \PosInt$
માટે $\{1, \dots, n\}$ છે.
\end{prop}

\begin{proof}
વિધેય $g$ને પુનરાવર્તી રીતે વ્યાખ્યાયિત કરીએ:
$g(1) = f(1)$ લો. જો $g(i)$ પહેલેથી વ્યાખ્યાયિત હોય,
તો $f(1)$, $f(2)$, \dots{}ની જે પહેલી કિંમત
$g(1)$, \dots, $g(i)$માં હજી આવી ન હોય,
તેને $g(i+1)$ લો, જો આવી કિંમત હોય.
જો $A$માં માત્ર $n$ !!{element}s હોય, તો
$g(1)$, \dots, $g(n)$ બધા વ્યાખ્યાયિત થાય છે,
એટલે આપણે વિધેય $g\colon \{1, \dots, n\}
\to A$ વ્યાખ્યાયિત કર્યું છે. જો $A$માં અનંત
સંખ્યામાં !!{element}s હોય, તો કોઈ પણ $i$ માટે
પરિગણના $f(1)$, $f(2)$, \dots{}માં $A$નો એવો
!!a{element} હોવો જ જોઈએ જે $g(1)$, \dots, $g(i)$માં
હજી આવ્યો ન હોય. આ કિસ્સામાં આપણે વિધેય
$g\colon \PosInt \to A$ વ્યાખ્યાયિત કર્યું છે.

વિધેય $g$ એ !!{surjective} છે, કારણ કે $A$નો
કોઈ પણ ઘટક $f(1)$, $f(2)$, \dots{}માં આવે છે
($f$ એ !!{surjective} હોવાથી), અને તેથી છેવટે
કોઈક~$i$ માટે $g(i)$ની કિંમત બનશે.
તે !!{injective} પણ છે: જો $g(j) = g(i)$ થાય
એવો $j < i$ હોત, તો $g(i)$ પહેલેથી જ
$g(1)$, \dots, $g(i-1)$માં હોત; તે $g$ની
આપણી વ્યાખ્યાથી વિરુદ્ધ છે.
\end{proof}

\begin{cor}\ollabel{cor:enum-nat-bij}
ગણ $A$ એ !!{enumerable} હોય જો અને માત્ર જો
તે ખાલી હોય અથવા !!a{bijection}
$f\colon N \to A$ હોય, જ્યાં $N = \Nat$
અથવા કોઈક $n \in \Nat$ માટે $N = \{0, \dots, n\}$ હોય.
\end{cor}

\begin{proof}
$A$ એ !!{enumerable} હોય જો અને માત્ર જો $A$
ખાલી હોય અથવા !!a{surjective} $f\colon \PosInt \to A$
હોય. \olref{prop:enum-bij} પ્રમાણે, બીજી શક્યતા
ત્યારે અને ત્યારે જ સાચી છે જ્યારે એવું !!a{bijective}
વિધેય~$f\colon Z \to A$ હોય કે $Z =
\PosInt$ અથવા કોઈક $n \in \PosInt$ માટે
$Z = \{1, \dots, n\}$ હોય.
\olref{prop:enum-shift}ની સાબિતી જેવી જ દલીલથી,
આ પણ ત્યારે અને ત્યારે જ થાય જ્યારે
!!a{bijection} $g\colon N \to A$ હોય, જ્યાં
$N = \Nat$ અથવા $N = \{0, \dots, n-1\}$ હોય.
\end{proof}

\begin{prob}
\olref[sfr][siz][enm]{defn:enumerable} પ્રમાણે ગણ $A$
ગણનીય હોય જો અને માત્ર જો $A = \emptyset$
હોય અથવા !!a{surjective} $f\colon
\PosInt \to A$ હોય. “!!{enumerable} ગણ”ની ચોક્કસ
વ્યાખ્યા આ રીતે પણ આપી શકાય: ગણ ગણનીય હોય
જો અને માત્ર જો !!a{injective} વિધેય
$g\colon A \to \PosInt$ હોય.
બતાવો કે વ્યાખ્યાઓ સમકક્ષ છે, એટલે કે
!!a{injective} વિધેય $g\colon A \to \PosInt$
હોય જો અને માત્ર જો $A = \emptyset$ હોય અથવા
!!a{surjective} $f\colon \PosInt \to A$ હોય.
\end{prob}

\end{document}
